\documentclass[12pt,addpoints,answers]{exam}
\usepackage[utf8]{inputenc}
\usepackage{amsmath,amsfonts,amssymb,amsthm}
\usepackage[margin=1in]{geometry}
\usepackage{mathtools}
\usepackage{hyperref}
\usepackage{fullpage}
\usepackage{microtype}
\usepackage{xspace}
\usepackage[svgnames]{xcolor}
\usepackage[sc]{mathpazo}
\usepackage{enumitem}
\usepackage{bm}

%----------Header--------------------%
\def\course{{\sc Introduccion a la Criptografia Moderna}}
\def\term{Depto. de Computaci\'{o}n, FCEN, UBA, 1er. Bimestre 2026}
\def\prof{Lecturer: Juan Garay}
\newcommand{\handout}[5]{
   \renewcommand{\thepage}{\arabic{page}}
   \begin{center}
   \framebox{
      \vbox{
    \hbox to 5.78in { \hfill \large{\course} \hfill }
    \vspace{2mm}
%    \hbox to 5.78in { \hfill \large{\prof} \hfill }
%       \vspace{2mm}
       \hbox to 5.78in { {\Large \hfill \textbf{#5}  \hfill} }
       \vspace{2mm}
       \hbox to 5.78in { \term \hfill \emph{#2}}
       \hbox to 5.78in { {#3 \hfill \emph{#4}}}
      }
   }
   \end{center}
   \vspace*{4mm}
}
\newcommand{\hw}[4]{\handout{#1}{#2}{#3}{#4}{Homework #1}}

% -- For ignoring stuff -- %
\newcommand{\ignore}[1]{}


\begin{document}

%----Specs: change accordingly-----%
\newif\ifstudent % comment out false
\studenttrue 
% \studentfalse

\def\hwnum{1}
\def\issuedate{26/3/26}
\def\duedate{7/4/25, 23:59 hs} % 
\def\yourname{} % put your name here
%------------------------------%

\ifstudent
\hw{\hwnum}{\issuedate}{Student: \yourname}{Due: \duedate}%
\else
\hw{\hwnum}{\issuedate}{\prof}{Due: \duedate}%
\fi

\noindent \textbf{Instructions}

\begin{itemize}
    \item Upload your solution to Campus; make sure it's only one file, and clearly write your name on the first page. Name the file \textsf{`$<$your last name$>$\_HW1.pdf.'} 
    %{\bf Important:} Make sure to tap {\bf Turn in} after you upload your solution.   
      
     If you are proficient with \LaTeX, you may also typeset your submission and submit in PDF format. To do so, uncomment the ``\%\textbackslash begin\{solution\}'' and ``\%\textbackslash end\{solution\}'' lines and write your solution between those two command lines.
    
      \item Your solutions will be graded on \emph{correctness} and
    \emph{clarity}. You should only submit work that you believe to be
    correct.
    % , and you will get significantly more partial credit if you     clearly identify the gap(s) in your solution.
    
    \item You may collaborate with others on this problem set.  However,
    you must \textbf{{write up your own solutions}} and \textbf{{list
      your collaborators and any external sources (including ChatGPT and similar generative AI chatbots)}} for each
    problem. Be ready to explain your solutions orally to a member of the course's staff
    if asked.
    
    \ignore{
    \item For problems that require you to provide an algorithm, you must
    give a precise description of the algorithm, together with a proof
    of correctness and an analysis of its running time. You may use
    algorithms from class as subroutines. You may also use any facts
    that we proved in class or from the book.
    } %IGNORE
    
\end{itemize}


\noindent This assignment contains \numquestions\ questions,
% \numpages\ pages
for a total of \numpoints \ points.
% and \numbonuspoints\ bonus points.

\bigskip
%\newpage

\begin{questions}


\question We saw the following definition of {\em negligible} functions in class:

{\bf Definition:} 
{\em A positive function $\mu(n)$ is called {\em negligible} if for any positive polynomial $p(n)$, there exists a constant $n_0$ such that for all $n \geq n_0$ we have $\mu(n) < 1/p(n)$.} 

\begin{parts}
    \part[6] Determine whether each of the following functions is negligible. Justify your answers.
    \begin{enumerate}[label=\roman*.]
        \item $f(n)=2^{-(\log n)(\log\log n)}$
        \begin{solution}
            Sin perdida de generalidad asumo que la base del logaritmo es 2, si no lo fuera queda una constante molestando en el exponente pero el argumento es el mismo. \\ \\
            $2^{-(\log n)(\log\log n)} = (2^{\log n})^{(-\log\log n)} = n^{((-\log\log n))} = \frac{1}{n^{((\log\log n))}}$. \\
            $\mu(n) = \frac{1}{n^{((\log\log n))}}$ es negligible porque: \\
            Sea $p(n)$ de grado $g$ tal que $p(n) = an^{g} + bn^{g-1} + \dots + c$ con $g > 0$ y las constantes perteneciente a los reales, siempre para un $n_0$ lo suficientemente grande $\log\log n_0 > g$, por lo que: $n^{\log\log n} > n^{g}$, por lo que $\frac{1}{n^{\log\log n}}< \frac{1}{n^{g}}$  
        \end{solution}
        \item $f(n)=n^{-1/n}$
        \begin{solution}
            $n^{-1/n} = \frac{1}{n^{1/n}}$. \\ \\
            $\mu(n) = \frac{1}{n^{1/n}}$ no es negligible porque: \\
            Como $\lim_{n \to \infty} \frac{1}{n} = 0$, se tiene que $\lim_{n \to \infty} n^{1/n} = 1$, por lo que \\ $\lim_{n \to \infty} \frac{1}{n^{1/n}} = 1$. \\
            Entonces, para $p(n) = n$, $\nexists\, n_0$ tal que $\forall\, n \geq n_0 \text{: } \frac{1}{n^{1/n}} < \frac{1}{n}$, porque $\lim_{n \to \infty} \frac{1}{n^{1/n}} = 1$, por lo tanto $\frac{1}{n^{1/n}} = 1 > \frac{1}{n}$ para un $n$ grande.
        \end{solution}
        \item $f(n)=1/n$ for odd $n$ and $f(n)=2^{-n}$ for even $n$
        \begin{solution}
            \[
            \mu(n) =
            \begin{cases}
            \frac{1}{n} & \text{si } n \text{ es impar},\\
            \frac{1}{2^n} & \text{si } n \text{ es par}.
            \end{cases}
            \]
            $\mu(n)$ no es negligible porque para $p(n) = n^{2}$ no se cumple que: $\exists \ n_0$ tal que $\forall n \geq n_0 \text{:} \ \mu (n) < 1/n^{2}$ porque para todos los $n$ impares $1/n > 1/n^{2}$
        \end{solution}
    \end{enumerate}
    \part[4] We saw that the definition can be restated as follows: 

    {\bf Lemma:} {\em A positive function $\mu(n)$ is {\em negligible} if and only if $\mu(n) \in O(1/p(n))$ for {\em all} positive polynomials $p$.}
    
    Prove the reverse ($\impliedby$) direction of this statement.
    \begin{solution}
        Que $\mu(n) \in O(1/p(n))$ $\forall$ polinomio positivo $p$ equivale por definición de Big O a decir que:
        $\exists c > 0,\ \exists n_0 \ \text{tal que} \ \forall n \ge n_0:\ \mu(n) \le \frac{c}{p(n)}$.

        Como vale para todo polinomio positivo, sea $h(n) = (c+1)p(n)$. Entonces,
        $\mu(n) \le \frac{c}{h(n)} = \frac{c}{(c+1)p(n)} < \frac{1}{p(n)}$.

        Por lo tanto, $\exists n_0$ tal que $\forall n \ge n_0:\ \mu(n) < \frac{1}{p(n)}$, y así $\mu(n)$ es negligible por definición.
    \end{solution}
\end{parts}

%\newpage

\question[10] Assume an attacker knows that a user’s password is either \texttt{abcd} or \texttt{bedg}.
Say the user encrypts his password using the shift cipher, and the attacker sees the resulting ciphertext. Show how the attacker can determine the user’s password, or explain why this is not possible.

\begin{solution}
    Con las posibles contraseñas no puede pasar que existan dos claves $k_1$ y $k_2$ tales que $Enc(\text{abcd}, k_1) = Enc(\text{bedg}, k_2)$ porque la diferencia entre las letras no es contante: las letras en posiciones impares de cada clave estan a distancia 1, pero las letras en posiciones pares están a distancia 3. Por lo que puede pasar que existan claves que nos den ciphers que son iguales en las posiciones pares (tomando $k_2 = k_1 + 3$) o en las impares (tomando $k_2 = k_1 + 1$), pero no en ambas. \\ \\ 
    Sabiendo que solo puede ser una de las dos contraseñas la que encripta al resultado, se me ocurren dos formas de verificar cual de las dos contraseñas es la correcta:
    \begin{itemize}
        \item La primer opción consiste en, por medio de fuerza bruta (ya que el espacio de claves es muy chico), probar para cada $k \in \{0, \dots , 25\}$ si se cumple:
        \[
            dec(\text{ciphertext}, k) = \text{abcd}
            \quad \text{o} \quad
            dec(\text{ciphertext}, k) = \text{bedg}.
        \]
        A la primera igualdad que valga, se determina cuál de las dos opciones es la contraseña.
        
        \item La segunda opción más inteligente es despejar $k_1$ y $k_2 \in \{ 0, \dots 25 \}$ tales que:
        \[
            1 + k_1 \equiv c_1 \pmod{26}
            \quad \text{y} \quad
            2 + k_2 \equiv c_1 \pmod{26},
        \]
        siendo $c_1$ el número asociado a la primera letra del ciphertext. Luego, comprobar si en la siguiente posición al sumar la letra por la clave se llega a la misma letra que el cipher. La clave que consiga las dos igualdades determina cuál de las dos contraseñas era la correcta.
    \end{itemize}
\end{solution}

%\newpage

\question[10] We proved in class that if a cipher is {\em perfectly secure/secret}, then it is {\em indistinguishable}. Prove the \textbf{opposite} direction. Begin by stating the perfect security and indistinguishability properties.
\begin{solution}
    Un cifrado es perfecto si y solo si:
    \[
        \forall m \in \mathcal{M}, \ \forall c:\quad 
        \Pr(\text{plaintext} = m \mid \text{ciphertext} = c) = \Pr(\text{plaintext} = m)
    \]
    Un cifrado es indistinguible si y solo si:
    \[
        \forall m_1, m_2 \in \mathcal{M}, \ \forall c:\quad 
        \Pr(\text{Enc}(m_1) = c) = \Pr(\text{Enc}(m_2) = c)
    \]
    Prueba: \\ 
    \[ 
        Pr(plaintext = m|ciphertext = c) = \frac{Pr(plaintext = m \land ciphertext = c)}{Pr(ciphertext = c)} = \\
        \frac{Pr(Enc(m) = c)}{Pr(ciphertext = c)}
    \]
\end{solution}

%\newpage

\question We look at the different conditions under which the shift and mono-alphabetic substitution ciphers are perfectly secret:
\begin{parts}
    \part[4] Prove that if only a single character is encrypted, then the shift cipher is perfectly secret.
    % \begin{solution}
    % % Uncomment and type your solution here
    % \end{solution}
    \part[3] Does the perfectly secret conditional probability equation hold for the mono-alphabetic substitution cipher? Elaborate.
    % \begin{solution}
    % % Uncomment and type your solution here
    % \end{solution}
    \part[3] What is the largest message space $\cal M$ for which the mono-alphabetic substitution cipher provides perfect secrecy?
    % \begin{solution}
    % % Uncomment and type your solution here
    % \end{solution}
\end{parts}

%\newpage

\question[10] Suppose you have a message $m$ of length $n$, but you can only generate random keys of length $k$ and $l$ such that $k + \ell = n - 1$. You decide to generate two random keys and combine them with an additional bit. Does the resulting encryption scheme $E_{k_1,k_2}$, defined as follows, provide perfect secrecy?
\begin{center}
      $E_{k_1,k_2}(m) = (k_1\mathbin\Vert 1 \mathbin\Vert k_2) \oplus m$,
\end{center}
where $m \in \{0,1\}^n$, $k_1 \in \{0,1\}^k$, $k_2 \in \{0,1\}^\ell$, and $\Vert$ denotes concatenation. What if instead, the message space is modified to only include messages for which the $(k+1)$th bit is $0$ (i.e., $\mathcal{M}=\big\{m_1\mathbin\Vert0\mathbin\Vert m_2\mathbin\big\vert m_1\in\{0,1\}^k,m_2\in\{0,1\}^\ell\big\}$)?
% \begin{solution}
% % Uncomment and type your solution here
% \end{solution}

%\newpage

\question [10] Assume we require only that an encryption scheme $(\mathsf{Gen}, \mathsf{Enc},\mathsf{Dec})$ with message space $\mathcal{M}$ satisfy the following: For all $m \in \mathcal{M}$, we have $Pr[\mathsf{Dec}_K(\mathsf{Enc}_K(m)) = m] \geq 2^{-t}$. (This probability is taken over the choice of the key as well as any randomness used during encryption.) Show that perfect secrecy can be achieved with $|\mathcal{K}| < |\mathcal{M}|$ when $t\geq 1$. 

% \begin{solution}
% % Uncomment and type your solution here
% \end{solution}

\end{questions}

\end{document}